%!TEX root = ../thesis.tex
\renewcommand{\abstractUkr}{\chapter*{РЕФЕРАТ}\thispagestyle{empty}\renewcommand{\arraystretch}{1.5}\setfontsize{14}}
\renewcommand{\abstractEng}{\chapter*{ABSTRACT}\thispagestyle{empty}\renewcommand{\arraystretch}{1.5}\setfontsize{14}}


\abstractUkr

Магістерську дисертацію виконано на \pageref{LastPage} аркушах, вона містить перелік посилань на використані джерела з 12 найменувань. У роботі наведено 1 рисунок та 1 таблицю.

\textbf{Актуальність теми.}
Автоматична ідентифікація музичних інструментів в аудіозаписах є
актуальною задачею музичної інформатики, яка знаходиться на перетині
цифрової обробки сигналів і машинного навчання. Визначення складу
інструментів у музичному записі є необхідним для автоматичної
каталогізації, пошуку та аналізу музичних колекцій, а також для
підтримки задач автоматичного аранжування, ремейксування та
реставрації аудіо.

Реальні музичні твори містять одночасне звучання кількох інструментів,
тому задача відноситься до класу multi-label класифікації, де кожному
аудіозапису може відповідати довільна підмножина інструментальних
міток. Застосування сучасних методів глибинного навчання, зокрема
згорткових та рекурентних нейронних мереж, відкриває можливості для
побудови ефективних автоматизованих систем ідентифікації інструментів
безпосередньо на основі аудіосигналу.

\textbf{Мета і задачі дослідження.}
Метою дисертаційної роботи є побудова математичного та програмного
забезпечення методу автоматичної ідентифікації музичних інструментів
у звуковому форматі та дослідження його ефективності для задачі
multi-label класифікації інструментального складу аудіозаписів.

Для досягнення поставленої мети вирішено такі задачі:

а) проведено огляд і порівняльний аналіз існуючих методів автоматичної
ідентифікації музичних інструментів в аудіозаписах;

б) виконано аналіз і порівняння підходів до побудови математичної
моделі методу multi-label класифікації інструментального складу,
обрано оптимальний варіант;

в) побудовано математичну модель методу, що описує повний цикл
перетворення аудіосигналу у вектор інструментальних міток;

г) проведено порівняльний аналіз варіантів конфігурації
запропонованого методу та досліджено його властивості на
стандартизованому наборі даних;

д) виконано порівняння ефективності запропонованого методу з
існуючими підходами.

\textbf{Об'єктом дослідження} є інструментальний склад музичних творів
у звуковому форматі та методи його автоматичної ідентифікації на основі
цифрової обробки сигналів і машинного навчання.

\textbf{Предметом дослідження} є математичне та програмне забезпечення
методу автоматичної ідентифікації музичних інструментів (гітара, фортепіано,
ударні, бас, скрипка тощо) на основі мел-спектрограм аудіосигналу та
гібридної архітектури CNN~+~BiLSTM.

\textbf{Методи дослідження.}
Методи цифрової обробки аудіосигналів (STFT, мел-фільтрбанк), методи
виділення тембральних спектральних ознак (мел-спектрограми, MFCC),
методи multi-label класифікації на основі згорткових і рекурентних
нейронних мереж, методи навчання та оптимізації моделей машинного
навчання, експериментальні методи порівняльного аналізу результатів.

\textbf{Ключові слова:}
ідентифікація музичних інструментів, multi-label класифікація,
машинне навчання, згорткові нейронні мережі, CNN, BiLSTM,
мел-спектрограма, MFCC, бінарна крос-ентропія, F1-міра, mAP,
музична інформатика, IRMAS, тембральний аналіз.

\abstractEng

The master's thesis comprises \pageref{LastPage} pages and includes a list of 12 references. The work contains 1 figure and 1 table.

\textbf{Relevance.}
Automatic identification of musical instruments in audio recordings is
a relevant problem in music information retrieval that lies at the
intersection of digital signal processing and machine learning.
Determining the instrument composition of a musical recording is
essential for automatic cataloguing, search, and analysis of music
libraries, as well as for supporting tasks such as automatic
arrangement, remixing, and audio restoration.

Real-world musical works contain the simultaneous sound of several
instruments, making this a multi-label classification problem where
each audio recording may correspond to an arbitrary subset of
instrument labels. The application of modern deep learning methods,
in particular convolutional and recurrent neural networks, opens
opportunities for building effective automated instrument identification
systems directly from the audio signal.

\textbf{Purpose and objectives.}
The aim of the dissertation is to construct the mathematical and software
support for a method of automatic identification of musical instruments
in audio format and to investigate its effectiveness for the multi-label
classification of the instrument composition of audio recordings.

To achieve this goal, the following tasks were solved:

a) a review and comparative analysis of existing methods for automatic
identification of musical instruments in audio recordings was conducted;

b) an analysis and comparison of approaches to constructing the
mathematical model of the multi-label instrument classification method
was performed, and the optimal variant was selected;

c) a mathematical model of the method was constructed, describing the
full pipeline from audio signal to an instrument label vector;

d) a comparative analysis of configuration variants of the proposed
method was conducted and its properties were investigated on a
standardized dataset;

e) the effectiveness of the proposed method was compared with existing
approaches.

\textbf{Object of research} is the instrument composition of musical
works in audio format and the methods of its automatic identification
based on digital signal processing and machine learning.

\textbf{Subject of research} is the mathematical and software support
for the method of automatic identification of musical instruments
(guitar, piano, drums, bass, violin, etc.) based on mel-spectrograms
of the audio signal and a hybrid CNN~+~BiLSTM architecture.

\textbf{Research methods.}
Methods of digital audio signal processing (STFT, mel filterbank),
methods of timbral spectral feature extraction (mel-spectrograms,
MFCC), methods of multi-label classification based on convolutional
and recurrent neural networks, methods of training and optimization
of machine learning models, experimental methods of comparative
analysis of results.

\textbf{Scientific novelty.}
For the first time, a method for automatic multi-label identification
of musical instruments in audio recordings has been proposed, based on
mel-spectrograms and a hybrid CNN~+~BiLSTM architecture. Unlike
existing approaches focused on single-label recognition or rule-based
filtering, the proposed method directly models the simultaneous
presence of multiple instruments by applying independent sigmoid
activations per class and optimizing binary cross-entropy loss.

\textbf{Practical value.}
Based on the proposed method, software has been implemented that
performs automatic identification of the instrument composition of
musical recordings in audio format. The method has been verified and
validated on the standardized IRMAS dataset, achieving F1-macro and
mAP scores that substantially exceed the results of baseline methods.
The obtained results may be used for the development of applied tools
for music search, automatic annotation, and content-based retrieval
in large audio collections.

\textbf{Keywords:}
musical instrument identification, multi-label classification,
machine learning, convolutional neural networks, CNN, BiLSTM,
mel-spectrogram, MFCC, binary cross-entropy, F1-score, mAP,
music information retrieval, IRMAS, timbral analysis.
\clearpage

\pagestyle{plain}
